% HW1 for Intro to Data Science
% also my first real foray into LaTeX ._.

\documentclass{article} 
\usepackage{fancyhdr, titling, graphicx, mathtools, fixltx2e, tikz}
\pagestyle{fancy}
\usetikzlibrary{positioning}
\newdimen\nodeDist
\nodeDist=35mm

% put name and page number in header
\fancyhead[L]{Green, Ben}
\fancyhead[R]{Page \thepage}

\author{Ben Green}
\title{Intro to Data Science HW 2}

\begin{document}

% title page
\begin{titlepage}
	\begin{center}
	\textsc{\LARGE Intro to Data Science HW 2}\\
	\vspace{3mm}
	
	{\large \theauthor}\\
	
	\tableofcontents
	\setcounter{secnumdepth}{0}
	\vfill
	
	{\large \today}
	\end{center}

\end{titlepage}

\section{Question 1}

\subsection{1a.}

\begin{equation}
n - m (1 - (\frac{m-1}{m})^n)
\end{equation}

\subsection{1b.}

\begin{equation}
1-(answer to 1a.)
\end{equation}

\section{Question 2}

\section{Question 5}

\begin{center}
\begin{tikzpicture}[
    node/.style={%
      draw,
      rectangle,
    },
  ]

    \node [node] (A) {A};
    \path (A) ++(-135:\nodeDist) node [node] (B) {yes};
    \path (A) ++(-45:\nodeDist) node [node] (C) {B};
    \path (C) ++(-135:\nodeDist) node [node] (D) {C};
    \path (C) ++(-45:\nodeDist) node [node] (E) {no};
    \path (D) ++(-135:\nodeDist) node [node] (F) {yes};
    \path (D) ++(-45:\nodeDist) node [node] (G) {no};

    \draw (A) -- (B) node [left,pos=0.25] {true}(A);
    \draw (A) -- (C) node [right,pos=0.25] {false}(A);
    \draw (C) -- (D) node [left,pos=0.25] {true}(A);
    \draw (C) -- (E) node [right,pos=0.25] {false}(A);
    \draw (D) -- (F) node [left,pos=0.25] {true}(A);
    \draw (D) -- (G) node [right,pos=0.25] {false}(A);
\end{tikzpicture}
\end{center}

\section{Question 7}

Assuming that you've picked door number 1, there are three (equally likely) possible scenarios:

\begin{enumerate}
\item You pick the door with the prize, and the other two doors are empty.
\item Both your door and door number 2 are empty.
\item Both your door and door number 3 are empty.
\end{enumerate}

Overall, there are two groups of possibilities - that you've picked a winning door (1/3) or you've picked an empty door (2/3).\\

If an empty door is revealed, it must either be door 2 or 3, because you picked door number 1. Now you are presented with the same two groups - but the second group's doors now have probabilities 0 and 2/3 of containing the prize, when originally they were 1/3 each.\\

Since your door (door number 1) has a probability of 1/3 and the other closed door has a probability of 2/3, you should switch doors.



\end{document}